\section{Kriptográfia – részletesen}

A kriptográfia két nagy csoportra bontható: szimmetrikus és aszimmetrikus algoritmusokra. A szimmetrikus titkosításnál ugyanaz a kulcs szükséges a titkosításhoz és visszafejtéshez (pl. AES), míg az aszimmetrikus (pl. RSA) különböző kulcsokat használ a publikus és privát műveletekhez.

\subsection{AES példák és magyarázat}
\begin{lstlisting}[language=bash]
# AES-128-CBC encryption with a password
openssl aes-128-cbc -k PASSWORD -in secret.txt -out secret.txt.aes -p
# Decrypt
openssl aes-128-cbc -d -k PASSWORD -in secret.txt.aes -out secret.txt.dec
\end{lstlisting}

Magyarázat: a \texttt{-k} egyszerű jelszót fogad; a valós környezetben \texttt{-pass} vagy kulcsfájl használata, PBKDF2 és só ajánlott: \texttt{-pbkdf2 -iter 200000}.

\subsection{RSA és kulcsmenedzsment}
A kulcsgenerálás, export és konverzió gyakori feladat; itt egy részletes példa a kulcsok kezelésére és a PKCS\#8 formátum használatára.
\begin{lstlisting}[language=bash]
# Generate RSA 2048 key pair (recommended minimum)
openssl genpkey -algorithm RSA -out key-2048.pem -pkeyopt rsa_keygen_bits:2048
# Extract public key
openssl rsa -in key-2048.pem -pubout -out key-2048.pub
# PKCS\#8 (unencrypted) conversion
openssl pkcs8 -topk8 -inform PEM -outform PEM -nocrypt -in key-2048.pem -out key-2048-pkcs8.pem
\end{lstlisting}

\subsection{Aláírások és ellenőrzések}
Digitális aláírás ellenőrizhető a \texttt{openssl dgst} parancsal; hasznos példák és leírások a cheatsheetben találhatók.
